% =============================================
% File: chapters/experiment_results.tex
% Chương 4: Kết quả thực nghiệm
% =============================================

\chapter{KẾT QUẢ THỰC NGHIỆM}
\label{chap:ketqua}

Chương này trình bày kết quả cụ thể của từng giai đoạn thực nghiệm đã được mô tả
trong Chương~\ref{chap:phuongphap}: hàm truyền động cơ nhận dạng được, bộ tham số
PI tối ưu, hiệu quả huấn luyện mạng ANN và kết quả đáp ứng hệ thống trên phần cứng
ESP32.


% ==============================================
\section{Kết quả nhận dạng hàm truyền động cơ NF5475}
% ==============================================

\subsection{Hàm truyền xác định được}

Dữ liệu thực nghiệm được đưa vào công cụ \texttt{tfest} trong MATLAB để nhận dạng
hàm truyền. Kết quả xác định mô hình động cơ:

\begin{equation}
    G(s) = \frac{353{,}5s^3 + 81{,}55s^2 + 16{,}27s + 3{,}984}
                {s^4 + 33s^3 + 7{,}444s^2 + 3489s + 0{,}2802}
    \label{eq:transfer_function}
\end{equation}

Đánh giá mô hình so với dữ liệu thực:
\begin{itemize}
    \item Độ chính xác: \textbf{95,36\%}
    \item FPE (Final Prediction Error): 2572
    \item MSE (Mean Squared Error): 2562
\end{itemize}


% ==============================================
\section{Kết quả tối ưu tham số PI bằng GA}
% ==============================================

\subsection{So sánh phương pháp tối ưu}

Ba phương pháp tối ưu được so sánh (chi tiết quy trình xem Chương~\ref{chap:phuongphap}):
Random gain in bounds, MATLAB PID Tuner, và Giải thuật di truyền (GA).
GA cho kết quả vượt trội với:
\begin{itemize}
    \item Độ vọt lố (overshoot): $< 5\%$
    \item Thời gian xác lập: $< 0{,}5$ s
    \item Độ khớp với dữ liệu (fit-to-data): $\approx 95\%$
\end{itemize}

\subsection{Bộ gain tối ưu}

Kết quả tối ưu hóa bằng GA tại setpoint 250 RPM:
\begin{itemize}
    \item $K_p^* = 0{,}185436$
    \item $K_i^* = 5{,}288458$
\end{itemize}

Quy trình được lặp lại cho 9 mức setpoint từ 250 RPM đến 2250 RPM (bước 250 RPM),
thu được 9 bộ gain tương ứng. Mỗi bộ gain được kiểm chứng trên dữ liệu thực
nghiệm 10 giây (1001 mẫu).

\subsection{Cài đặt bộ PI trên ESP32}

Bộ PI vận tốc được cài đặt theo công thức rời rạc với $\Delta t = 0{,}01$ s:

\begin{align}
    e_{vel}(k) &= v_{req}(k) - v_{cur}(k) \\
    \int e\,dt &\mathrel{+}= e_{vel}(k) \cdot \Delta t \\
    u(k) &= K_p \cdot e_{vel}(k) + K_i \cdot \int e\,dt
\end{align}

Anti-windup: giới hạn tích phân trong $[-1000, 1000]$.
Giới hạn PWM đầu ra: $u(k) \in [-255, 255]$.


% ==============================================
\section{Kết quả huấn luyện mạng ANN}
% ==============================================

\subsection{Tóm tắt kiến trúc và cấu hình}

Kiến trúc và quy trình huấn luyện được trình bày chi tiết trong
Chương~\ref{chap:phuongphap}. Tóm tắt:
\begin{itemize}
    \item Kiến trúc: $2 \to 8 \to 16 \to 2$ (Dense + L2 + ReLU + Dropout).
    \item Optimizer: Adam, lr $= 2\times10^{-3}$, decay $= 3\times10^{-4}$.
    \item Số epochs: 1\,000.
\end{itemize}

\subsection{Kết quả huấn luyện}

Mạng ANN hội tụ sau khoảng 400 epochs. Đường cong loss huấn luyện và validation
bám sát nhau, chứng tỏ mạng không bị overfitting nhờ L2 regularization và Dropout
(xem Hình~\ref{fig:loss_curve}).

\begin{figure}[H]
    \centering
    % \includegraphics[width=0.75\textwidth]{figure/training_loss.png}
    \caption{Đường cong Training Loss và Validation Loss theo epoch}
    \label{fig:loss_curve}
\end{figure}

\subsection{Đánh giá mô hình}

Mô hình được đánh giá trên tập test (2 bộ dữ liệu chưa dùng để huấn luyện).
Kết quả dự đoán $K_p$, $K_i$ so với giá trị tối ưu từ GA cho thấy sai số nhỏ
và ổn định trên toàn dải setpoint 250 -- 2250 RPM.

\textit{Hạn chế hiện tại:} Chưa có số liệu định lượng đầy đủ về độ chính xác dự đoán
trên phần cứng thực; đây là hướng cần bổ sung trong giai đoạn tiếp theo.


% ==============================================
\section{Kết quả thực nghiệm trên phần cứng ESP32}
% ==============================================

\subsection{Kết quả thực nghiệm trên phần cứng}

Hệ thống PI vận tốc hoạt động ổn định trên ESP32 với 9 mức setpoint từ 250 đến
2250 RPM. Đáp ứng vận tốc đạt trong dải $\pm 5\%$ setpoint sau khoảng
$t_{rise} < 0{,}5$ s ở hầu hết các mức tải thử nghiệm
(xem Hình~\ref{fig:velocity_response}).

\begin{figure}[H]
    \centering
    % \includegraphics[width=0.85\textwidth]{figure/velocity_response.png}
    \caption{Đáp ứng vận tốc tại 9 mức setpoint (250 -- 2250 RPM)}
    \label{fig:velocity_response}
\end{figure}
